\documentclass[10pt,a4paper]{amsart}
\usepackage[margin=19mm]{geometry}
\usepackage{amsmath,amssymb,mathtools}
\usepackage[scr=rsfs,cal=euler]{mathalfa}
\usepackage{xfrac}
\usepackage[unicode,hidelinks,bookmarks=false]{hyperref}
\newcommand{\ie}{i.e.\ }
\newcommand{\cf}{cf.\ }
\newcommand{\resp}{respectively\ }
\newcommand{\Subsection}{\subsection}
\providecommand{\nobreakdash}{\nobreak\hbox{-}}
\setlength{\emergencystretch}{2em}
\multlinegap0pt
\usepackage{mathtools}
\usepackage{stackengine}
\usepackage{enumerate}
\usepackage{amssymb}
\newtheorem{theorem}{Theorem}
\newtheorem{lemma}{Lemma}
\newcommand{\D}{{\sf {D}}}
\def\Im{\protect\operatorname{Im}}
\def\Q{{\mathbb Q}}
\DeclareMathOperator{\cat}{{\mathsf{cat}}}
\DeclareMathOperator{\cu}{{\mathsf{cup}}}
\newcommand{\dD}{{\sf {dD}}}
\DeclareMathOperator{\dcat}{{\mathsf{dcat}}}
\def\id{\protect\operatorname{Id}}
\newcommand{\sd}{{(\sf {d})}}
\title{Homotopic distances and group-like spaces}
\author{Ekansh Jauhari, John Oprea}
\date{Source: arXiv:2607.03484v1 (CC BY 4.0)}
\begin{document}
\maketitle

\noindent\emph{Source and licence.} Homotopic distances and group-like spaces. Version arXiv:2607.03484v1 (CC BY 4.0), distributed by arXiv under the Creative Commons Attribution 4.0 International licence, http://creativecommons.org/licenses/by/4.0/, as recorded in the arXiv licence metadata for this exact version and shown on the abstract page https://arxiv.org/abs/2607.03484v1. Full licence notice: \texttt{licenses/arxiv-2607.03484-CC-BY-4.0.txt}.\par\smallskip

\noindent\textit{Editorial scope.} Counted unit: the article's theorem thm:nilfgstar (source0.tex lines 1761--1777). The article's complete original proof of it is delivered with it (source0.tex lines 1822--1872, one \texttt{proof} environment of 51 lines, which the article places away from the statement). No mathematics is written, repaired or re-derived: the statement and the proof are the article's own lines, in the article's own notation, and no retained line is substituted or re-flowed.  Retained uncounted as the source-local context the proof needs: lines 1226--1228; lines 1233--1240; lines 1808--1819.  Nothing was omitted from any retained span, so the package carries no omission record.  No retained line contains a citation, so the wrapper carries no bibliography and no bibliography entry.  No retained line contains a figure environment, an image-inclusion command or drawing code, so no image is delivered and the package declares no asset dependency.  Editorial formatting only, in addition: an AMS \texttt{amsart} preamble that restates the article's own declarations and macros used by the retained lines, namely the environments \texttt{theorem}, \texttt{lemma} on separate counters, together with the standard packages those lines need.  The retained environments print in this document as Theorem 1, Lemma 1 in order of appearance, whereas the article prints the same items with the numbering of its own full document.  The complete native source of the article as distributed in the arXiv e-print for this exact version is bundled unchanged as \texttt{sources/204392/source0.tex}.
\begin{equation}\label{eq: map Fm}
F_m(x_1,x_2,\dots,x_m)=(F(x_1,x_2),F(x_2,x_3),\dots, F(x_{m-1},x_m)).
\end{equation}

\begin{theorem}\label{thm: main1}
    Let $G$ be a path-connected group-like space or a path-connected CW $H$-space, and $m\ge 2$ be an integer. If $f_i\colon X\to G$ are maps for $1\le i \le m$, then 
    \[
    \sd\D(f_1,\dots,f_m)=\sd\cat(F_m\circ (f_1,\dots,f_m)),
    \]
    where $(f_1,\dots,f_m)\colon X\to G^m$ is the map $(f_1,\dots,f_m)(x)=
(f_1(x),\dots,f_m(x))$, and $F_m\colon G^m\to G^{m-1}$ is as in~\eqref{eq: map Fm}.
\end{theorem}

\begin{lemma}\label{lem:catf}
Suppose $X$ is a simply connected finite CW complex and $G$ is a compact group-like 
space. Let $g \in [X,G]$, and let its rationalization $g_\Q \colon X_\Q \to G_\Q$ be represented 
as above by $g_\Q=(g_1,\ldots,g_p)$, where $G\simeq_\Q \prod_{k=1}^p 
K(\Q,n_k)$ and 
$g_k = g_\Q^*(\beta_k) \in H^{n_k}(X;\Q)$ for a generator 
$\beta_k \in H^{n_k}(K(\Q,n_k);\Q)$. Then
\[
\cat(g_\Q)=\cu_\Q(\textup{Im}(g^*))=\cu_\Q\{g_1,\ldots,g_p\}, 
\]
where $\cu_\Q\{g_1,\ldots,g_p\}$ is the length of the longest non-zero cup product(s) of the form $g_{j_1}\smile \cdots\smile g_{j_t}\ne 0$ and $j_k\in\{1,\dots,p\}$ for $1\le k \le t$.  If, additionally, $G$ has the homotopy type of a CW complex, then $\dcat(g_\Q)=\cu_\Q\{g_1,\ldots,g_p\}$ as well.
\end{lemma}

\begin{theorem}\label{thm:nilfgstar}
Suppose $G$ is a path-connected compact group-like space of finite type and $X$ is a 
simply connected finite CW complex. For maps $f_i\colon X \to G$, $1\le i \le m$, denote
the rationalizations by $f_{i\Q}\colon X_\Q \to G_\Q$. Then we have that
\[
\D(f_{1\Q},\dots,f_{m\Q})=\cu_{\Q}\left(\Im\left(
%f_1^*-f_m^*
%%\bigotimes_{i=1}^{m-1}\left(f_i^*-f_{i+1}^*\right) 
%%%\bigl(f_1^*-f_{2}^*\bigr),\bigl(f_2^*-f_3^*\bigr),\dots,\left(f_{m-1}^*-f_{m}^*\right)
\Delta_{m-1}^*\circ \bigotimes_{i=1}^{m-1}\left(f_i^*-f_{i+1}^*\right)\right)
\right),
\]
where 
$f_i^*$ is the map induced in rational cohomology by $f_i$ and 
$\Delta_{m-1}\colon X\to X^{m-1}$ is the diagonal map. If, additionally, $G$ has the 
homotopy type of a CW complex, then $\dD(f_{1\Q},\dots,f_{m\Q})=\D(f_{1\Q},\dots,f_{m\Q})$.
\end{theorem}

\begin{proof}[Proof of Theorem~\ref{thm:nilfgstar}]
For convenience, we write $f_i$ instead of the more unwieldy $f_{i\Q}$, but the
reader should always have the rationalized maps in mind.

First, we have the equality $\sd\D(f_1,\dots,f_m)=\sd\cat(F_m\circ (f_1,\dots,f_m))$ from Theorem~\ref{thm: main1}, 
where $F_m\colon G^m\to G^{m-1}$ is defined in~\eqref{eq: map Fm} using the composition $F=\mu\circ(\id_G\times\eta)\colon G\times G\to G$ for a multiplication $\mu$ and an 
inversion $\eta$ on $G$. Taking $F_m\circ (f_1,\dots,f_m)\in [X,G^{m-1}]_\Q$ in 
Lemma~\ref{lem:catf}, we deduce that
\[
\D(f_1,\dots,f_m)=\cu_\Q\left(\Im\left(F_m\circ (f_1,\dots,f_m)\right)^*\right).
\]
Indeed, $G^{m-1}$ is compact and group-like, so everything said in this section so far 
for $G$ applies to $G^{m-1}$ as well. 
Furthermore, if $G$ has the homotopy type of a CW complex, then $G^{m-1}$ has that of a 
finite CW complex, so that we also have
\[
\dD(f_1,\dots,f_m)=\cu_\Q\left(\Im\left(F_m\circ (f_1,\dots,f_m)\right)^*\right).
\]
For convenience, let us write $F_m\circ (f_1,\dots,f_m)=(f_1-f_2,f_2-f_3,\dots,f_{m-1}-f_m)$, where for each $i\in\{1,\dots,m-1\}$, the map $f_i-f_{i+1}\colon X_\Q\to G_\Q$ is the composition
\[
X _\Q \xrightarrow{\Delta}X _\Q\times X _\Q\xrightarrow{f_i\times f_{i+1}} G _\Q\times 
G _\Q\xrightarrow{1\times\eta} G _\Q\times G_\Q \xrightarrow{\mu} G _\Q,
\]
so that the map $F_m\circ (f_1,\dots,f_m)$
%$=(f_1-f_2,f_2-f_3,\dots,f_{m-1}-f_m)$ 
can also be seen simply as the composition
\begin{equation}\label{eq: writing Fm composition}
X _\Q \xrightarrow{\Delta_{m-1}} X^{m-1} _\Q \xrightarrow{\prod_{i=1}^{m-1}(f_i-f_{i+1})} 
G^{m-1} _\Q.
\end{equation}
Now, $H^*(G;\Q)\cong\bigwedge(\beta_{n_k})$, an exterior algebra on the generators 
of $H^{n_k}(K(\Q,n_k);\Q)$ for $1\le k\le p$, and it is a well-known fact that these generators are primitive, i.e.,
\[
\mu^*(\beta_{n_k}) = \beta_{n_k} \otimes 1 + 1 \otimes \beta_{n_k}.
\]
We see for the induced map $(f_i-f_{i+1})^*\colon H^*(G;\Q)\to H^*(X;\Q)$ that 
\begin{align*}
(f_i-f_{i+1})^*(\beta_{n_k}) & = \Delta^*((f_i^*\times f_{i+1}^*)((1\times \eta^*)(\mu^*(\beta_{n_k})))) \\
& = \Delta^*((f_i^*\times f_{i+1}^*)((1\times \eta^*)(\beta_{n_k} \otimes 1 + 1 \otimes \beta_{n_k}))) \\
& =  \Delta^*((f_i^*\times f_{i+1}^*)(\beta_{n_k} \otimes 1 - 1 \otimes \beta_{n_k})) \\
& = \Delta^*(f_i^*(\beta_{n_k}) \otimes 1 - 1 \otimes f_{i+1}^*(\beta_{n_k})) \\
& = f_i^*(\beta_{n_k}) - f_{i+1}^*(\beta_{n_k}),
\end{align*}
since $\Delta^*\colon H^*(X\times X;\Q)\to H^*(X;\Q)$ gives the cup product. 
%Because the $\beta_{n_k}$ generate $H^*(G;\Q)$, we see that every element in $\Im((f_i-f_{i+1})^*)$ is a product of elements of the form $f_i^*(\beta_{n_k}) -f_{i+1}^*(\beta_{n_k})$.
Similarly, we see from~\eqref{eq: writing Fm composition} that $(F_m\circ (f_1,\dots,f_m))^*\colon H^*(G^{m-1};\Q)\cong H^*(G;\Q)^{\otimes m-1}\to H^*(X;\Q)$ on a generator $\beta_{1,n_k}\otimes\cdots\otimes \beta_{m-1,n_k}\in H^*(G;\Q)^{\otimes m-1}$ is
\begin{align*}
(F_m\circ (f_1,\dots,f_m))^*(\beta_{1,n_k}\otimes\cdots\otimes \beta_{m-1,n_k}) & = \Delta_{m-1}^*\left(\bigotimes_{i=1}^{m-1}\left(f_i^*\bigl(\beta_{i,n_k}\bigr) - f_{i+1}^*\bigl(\beta_{i,n_k}\bigr)\right)\right).
\end{align*}
Because the elements $\beta_{n_k}$ generate $H^*(G;\Q)$, and hence $\beta_{1,n_k}\otimes\cdots\otimes \beta_{m-1,n_k}$ generate $H^*(G^{m-1};\Q)$, we see that every element in $\Im(F_m\circ (f_1,\dots,f_{m}))^*$ is a product of elements as on the right side above. This completes the proof.
\end{proof}

\end{document}
